\section{Executive Brief: DUALEXIS Pilot Deployment}
\label{sec:executive-summary}

\subsection{Purpose}

DUALEXIS is a \textbf{governed pilot approach} to privacy-preserving \textbf{situational awareness}
and \textbf{emergency coordination} in schools.
This brief is for institutional decision-makers: leadership, safety coordinators, DPOs, ethics
committees, and Horizon consortium partners.
It supports planning conversations; it is not legal advice, certification, or proof of safety outcomes.

\subsection{What the pilot helps schools evaluate---not promise}

Licensed safety staff would receive \textbf{zone-level summaries} (e.g., crowding near an exit,
unusual activity in a cafeteria \emph{zone}) and \textbf{advisory} text that always passes
\textbf{human review} before any institutional action.

The pilot tests whether this \emph{supports} existing emergency plans---it does not replace them.

Processing stays on \textbf{school-controlled systems}.
Video and audio are \textbf{not retained by default}; facial recognition, biometrics, and identity
profiling are excluded (Box~\ref{box:system-boundaries}).

\subsection{Document pack and how to read it}

This \textbf{Executive Brief} is the entry point for non-technical stakeholders.
The companion \textbf{Deployment and Governance Requirements} document (\texttt{main.pdf},
v\docversion) contains:
\begin{itemize}[noitemsep]
  \item deployment and network rules;
  \item the full 15-item DPIA preparation checklist;
  \item standard operating procedures and mandatory tabletop exercises;
  \item pilot phases with go/no-go gates and stop conditions;
  \item governance templates (charter, transparency notice, protocol matrix, workshop agendas);
  \item institutional circulation and workshop guide (community communication, objection handling);
  \item indicative budgets, \textbf{Appendix~E hardware BOM}, and the full risk register.
\end{itemize}

\subsection{Why consider a governed pilot}

Schools need timely awareness during stressful events while communities resist identity-centric
monitoring.
A careful, co-designed pilot can:
\begin{itemize}[noitemsep]
  \item test whether zone summaries help staff build shared awareness \textbf{without naming individuals};
  \item rehearse alignment with existing lockdown, evacuation, and verification procedures;
  \item document accountability through review queues and audit trails;
  \item support Horizon or university research \textbf{only under explicit school governance}.
\end{itemize}

\textbf{Success now} means pilot \emph{readiness}---authorization, DPIA progress, protocol mapping,
training, ethics approval---not proof that technology alone prevents harm.

\subsection{Recommended pathway}

\begin{table}[htbp]
  \centering
  \caption{Indicative pilot pathway.}
  \label{tab:pilot-tiers-summary}
  \small
  \begin{tabular}{@{}p{2.4cm}p{2.8cm}p{5.8cm}@{}}
    \toprule
    Phase & Scale & Focus \\
    \midrule
    Foundation & 1 school, few zones &
      Charter, transparency, protocol mapping, table-top exercises. \\
    Recommended & 1--3 schools &
      Staff training; bounded live sensing (staffed hours). \\
    Extended & 5--10 schools &
      Multi-site learning; consortium replication (e.g.\ Horizon). \\
    \bottomrule
  \end{tabular}
\end{table}

\subsection{Co-design, opt-out, and phased participation}

\begin{itemize}[noitemsep]
  \item \textbf{Co-design:} schools choose which zones are in scope with safety staff and leadership;
        sensitive areas require explicit approval.
  \item \textbf{Phased participation:} live sensing only after sandbox training and \textbf{two tabletop drills}.
  \item \textbf{Opt-out:} zones can be excluded or sensing paused; stakeholder withdrawal triggers stop conditions.
  \item \textbf{Community voice:} parent councils or ethics boards consulted per local law before Phase~2.
\end{itemize}

\subsection{Human-in-the-loop and advisory-only operation}

Every elevated summary passes \textbf{human review} before any follow-up.
Staff may \emph{monitor}, \emph{verify}, or \emph{consider escalation}---but the system does not
execute lockdowns, call emergency services, or contact parents.
Licensed safety personnel and leadership retain full authority.
Figure~\ref{fig:governance-hitl-flow} shows the institutional flow.

\subsection{Governance readiness before live sensing}

\begin{enumerate}[noitemsep]
  \item leadership authorizes scope and community communication;
  \item site-specific DPIA with qualified counsel; ethics sign-off if evaluating staff;
  \item semantic events mapped to the school emergency plan;
  \item staff trained to treat outputs as \emph{advisory}, not orders;
  \item retention, access roles, false-alarm review, and \textbf{pause/stop} criteria agreed.
\end{enumerate}

\subsection{Go / no-go and stop (summary)}

\textbf{Go:} signed charter; DPIA acceptance; protocol matrix; $\geq$90\% training; two table-top exercises;
IT sign-off; ethics approval if applicable.

\textbf{Stop:} privacy violation; unauthorized access; regulatory inquiry; excessive false alarms;
attempted autonomous actuation.

Full detail: requirements document Sections~\ref{sec:pilot-programme} and~\ref{sec:operational-protocols}.

\noindent\textbf{Box 1 --- System boundaries, exclusions, and advisory role}
\label{box:system-boundaries}
\label{sec:what-dualxis-is-not}

\vspace{0.35em}
\noindent\fbox{\parbox{0.97\textwidth}{\small
\textbf{DUALEXIS is \emph{not}:}
\begin{itemize}[noitemsep,leftmargin=*]
  \item facial recognition, biometrics, or voiceprint identification;
  \item predictive policing, behavioral profiling, or automated disciplinary scoring;
  \item autonomous lockdown, alarm actuation, or emergency-service dispatch;
  \item cloud-based raw video/audio surveillance by default;
  \item student or staff tracking across zones, days, or campuses;
  \item identity-linked operational outputs (student IDs, staff directories).
\end{itemize}
\textbf{Non-negotiable pilot properties:}
\begin{itemize}[noitemsep,leftmargin=*]
  \item \textbf{Advisory only} --- supports staff comprehension; does not issue orders or
        trigger protocols automatically.
  \item \textbf{Human authority} --- PA, radio, parent communication, and escalation remain
        school responsibilities.
  \item \textbf{Privacy by design} --- zone-level summaries; no default media archive;
        local school-controlled processing.
  \item \textbf{Human-in-the-loop} --- elevated items require staff review before external action.
\end{itemize}
\textbf{The pilot provides} plain-language \textbf{zone summaries} for trained staff (e.g., crowding
near an exit, unusual activity in a cafeteria zone) plus \textbf{advisory text} that staff may
accept, modify, or dismiss.
Exclusions are fixed design rules---not optional settings without DPIA update, ethics review, and
steering approval.
}}
\vspace{0.7em}


% Institutional diagram: human-in-the-loop governance flow
\begin{figure}[htbp]
  \centering
  \begin{tikzpicture}[
    node distance=0.55cm and 0.45cm,
    box/.style={
      draw=black!75,
      rounded corners=2pt,
      minimum height=0.85cm,
      minimum width=2.35cm,
      align=center,
      font=\small,
      fill=white,
    },
    proc/.style={box, fill=black!4},
    human/.style={box, fill=black!8, line width=0.9pt},
    forbid/.style={
      box,
      dashed,
      draw=black!45,
      font=\footnotesize,
      minimum width=2.8cm,
    },
    lbl/.style={font=\footnotesize, text=black!65},
    arr/.style={-{Stealth[length=2.2mm]}, thick, draw=black!70},
  ]
    \node[proc] (sens) {Approved\\sensing};
    \node[proc, right=of sens] (sem) {Zone-level\\semantic events};
    \node[human, right=of sem] (rev) {Human\\review};
    \node[human, right=of rev] (dec) {Institutional\\decision};
    \node[proc, right=of dec] (prot) {Existing\\emergency protocols};

    \draw[arr] (sens) -- (sem);
    \draw[arr] (sem) -- (rev);
    \draw[arr] (rev) -- (dec);
    \draw[arr] (dec) -- (prot);

    \node[lbl, above=0.15cm of sem] {advisory inputs only};
    \node[lbl, above=0.15cm of rev] {mandatory authority};
    \node[forbid, below=1.05cm of rev] (noauto) {No autonomous lockdown, dispatch,\\or disciplinary actuation};
    \draw[dashed, draw=black!35] (rev.south) -- (noauto.north);
  \end{tikzpicture}
  \caption{Human-in-the-loop governance flow. DUALEXIS supplies structured summaries and
  \emph{advisory} recommendations; licensed personnel and school leadership retain
  operational authority. Emergency protocols are executed outside the system.}
  \label{fig:governance-hitl-flow}
\end{figure}


% Institutional diagram: governance and oversight model
\begin{figure}[htbp]
  \centering
  \begin{tikzpicture}[
    centerbox/.style={
      draw=black!80,
      line width=0.9pt,
      rounded corners=3pt,
      minimum width=3.4cm,
      minimum height=1.1cm,
      align=center,
      font=\small\bfseries,
      fill=black!6,
    },
    role/.style={
      draw=black!70,
      rounded corners=2pt,
      minimum width=2.85cm,
      minimum height=0.72cm,
      align=center,
      font=\footnotesize,
      fill=white,
    },
    arr/.style={-{Stealth[length=1.8mm]}, draw=black!50},
  ]
    \node[centerbox] (steer) {Pilot steering\\committee};

    \node[role, above left=1.1cm and 0.2cm of steer] (lead) {School\\leadership};
    \node[role, above=1.15cm of steer] (safe) {Safety\\officers};
    \node[role, above right=1.1cm and 0.2cm of steer] (dpo) {Data protection\\officer (DPO)};

    \node[role, below left=1.1cm and 0.2cm of steer] (eth) {Ethics\\committee};
    \node[role, below=1.15cm of steer] (it) {IT / network\\team};
    \node[role, below right=1.1cm and 0.2cm of steer] (res) {Research team\\(UCJC / partners)};

    \foreach \r in {lead,safe,dpo,eth,it,res} {
      \draw[arr] (\r) -- (steer);
    }

    \node[font=\footnotesize, text=black!55, align=center, text width=10cm]
      at ([yshift=-2.05cm]steer) {%
        Accountable decisions, DPIA, transparency, protocol mapping, and go/no-go for live sensing
        require alignment across roles---no single function owns the full pilot alone.};
  \end{tikzpicture}
  \caption{Governance and oversight model for a governed school pilot.}
  \label{fig:governance-oversight-model}
\end{figure}


% Institutional diagram: pilot phase roadmap with governance gates
\begin{figure}[htbp]
  \centering
  \begin{tikzpicture}[
    phase/.style={
      draw=black!75,
      rounded corners=2pt,
      minimum width=3.1cm,
      minimum height=0.7cm,
      align=center,
      font=\small,
      fill=black!5,
    },
    gate/.style={
      draw=black!55,
      rounded corners=1pt,
      minimum width=2.9cm,
      minimum height=0.55cm,
      align=center,
      font=\footnotesize,
      fill=white,
    },
    arr/.style={-{Stealth[length=2mm]}, thick, draw=black!70},
    node distance=0.38cm,
  ]
    \node[phase] (p0) {Phase 0\\Foundation};
    \node[gate, below=of p0] (g0a) {Leadership authorization};
    \node[gate, below=of g0a] (g0b) {Transparency plan};

    \node[phase, right=1.35cm of p0] (p1) {Sandbox /\\Tabletop};
    \node[gate, below=of p1] (g1a) {Protocol mapping};
    \node[gate, below=of g1a] (g1b) {Staff training (awareness)};

    \node[phase, right=1.35cm of p1] (p2) {MVP pilot};
    \node[gate, below=of p2] (g2a) {DPIA initiated or completed};
    \node[gate, below=of g2a] (g2b) {Ethics approval (if required)};

    \node[phase, right=1.35cm of p2] (p3) {Recommended\\pilot};
    \node[gate, below=of p3] (g3a) {Bounded live sensing};
    \node[gate, below=of g3a] (g3b) {Calibration review};

    \node[phase, right=1.35cm of p3] (p4) {Extended\\programme};
    \node[gate, below=of p4] (g4a) {Multi-site governance};
    \node[gate, below=of g4a] (g4b) {Consortium reporting (e.g.\ Horizon)};

    \draw[arr] (p0.east) -- (p1.west);
    \draw[arr] (p1.east) -- (p2.west);
    \draw[arr] (p2.east) -- (p3.west);
    \draw[arr] (p3.east) -- (p4.west);

    \node[font=\footnotesize\itshape, text=black!55, below=0.95cm of g0b]
      {Advance only when gates for the current phase are documented};
  \end{tikzpicture}
  \caption{Pilot phase roadmap with indicative governance gates. Timelines are
  institution-specific; no phase should add live occupant sensing without prior approval.}
  \label{fig:pilot-phase-roadmap}
\end{figure}


\subsection{What this brief does not claim}

\begin{itemize}[noitemsep]
  \item GDPR or EU AI Act certification;
  \item evidence that the pilot reduces harm or improves emergency outcomes;
  \item permission to process personal data without institutional authorization;
  \item binding budget commitments (see Appendix~B in the full requirements pack).
\end{itemize}

\subsection{Privacy, GDPR, and AI Act (planning summary)}

Schools must complete a site-specific \textbf{data protection impact assessment (DPIA)} with
qualified counsel.
Processing likely involves children and systematic monitoring in public areas of the school.
The reference design minimizes data: zone summaries, no biometrics, no default video archive,
local control.

\textbf{GDPR planning topics} include lawful basis, minimization, purpose limitation, retention,
accountability, and data-subject rights via the school DPO.
\textbf{EU AI Act planning} may apply to staff-facing decision support; human oversight and
transparency are required---but this brief does \textbf{not} certify conformity.

Full 15-item DPIA checklist: companion requirements document (\texttt{main.pdf}, Privacy and DPIA section).

\subsection{Key risks leadership should discuss}

\begin{table}[htbp]
  \centering
  \caption{Priority risks (summary).}
  \label{tab:exec-risks}
  \small
  \begin{tabular}{@{}clp{6.5cm}@{}}
    \toprule
    ID & Level & Theme \\
    \midrule
    R1 & High & Misread as surveillance $\rightarrow$ transparency, training \\
    R4 & High & Accidental media retention $\rightarrow$ audits, DPIA controls \\
    R6--R7 & High & Scope creep / autonomous actuation $\rightarrow$ contracts, SOP \\
    R10 & High & Children's data $\rightarrow$ DPO oversight \\
    R15 & Med & Community opposition $\rightarrow$ co-design, opt-out \\
    \bottomrule
  \end{tabular}
\end{table}

Full register: Appendix~C in \texttt{main.pdf}.

\subsection{Questions stakeholders often ask}

\paragraph{Is this facial recognition?}
No. The design excludes facial recognition, biometrics, and identity-linked outputs.

\paragraph{Does the system lock down the school automatically?}
No. Outputs are \textbf{advisory only}. Staff and leadership decide all protocol steps.

\paragraph{Where is video stored?}
Raw video and audio are not retained by default on DUALEXIS processing paths.
Only short-lived buffers and structured zone summaries apply, as defined in the site DPIA.

\paragraph{Can we stop or shrink the pilot?}
Yes. Zones can be excluded, sensing paused, or the pilot halted under documented stop conditions.

\paragraph{What must happen before cameras go live?}
Governance charter, DPIA acceptance, protocol mapping, staff training, IT sign-off, and
\textbf{two tabletop exercises} with remediations documented.

\paragraph{Will researchers control our emergency response?}
No. Schools retain operational authority. Research partners support training and evaluation
under steering-approved boundaries only.

\paragraph{Does this document certify GDPR compliance?}
No. It supports DPIA preparation and institutional planning; legal sign-off remains with the school.


\subsection{Horizon Europe and research partners}

Consortium partners should harmonize \textbf{governance templates}, summary formats, and
\textbf{anonymized} reporting boundaries---not operational control.
Research export remains \textbf{disabled by default}; schools authorize any learning data leaving site.

\subsection{Institutional boundaries at a glance}

\begin{center}
\fbox{\parbox{0.92\textwidth}{\small
\textbf{Advisory only.} The system supports staff comprehension; it does not issue orders,
trigger lockdowns, contact emergency services, or discipline students.\\
\textbf{Human authority.} Protocol activation, PA/radio use, parent communication, and
escalation remain school responsibilities.\\
\textbf{Privacy by design.} No biometrics; no default media archive; local school control.\\
\textbf{Tabletop obligation.} At least two exercises with synthetic or sandboxed feeds before
any live occupant sensing.
}}
\end{center}

\subsection{Who should be involved}

\begin{itemize}[noitemsep]
  \item \textbf{School leadership} --- scope, charter, community communication;
  \item \textbf{Safety coordinators} --- protocol mapping, review workflows, tabletops;
  \item \textbf{DPO} --- DPIA, lawful basis, transparency, subject rights;
  \item \textbf{Ethics committee} --- human-subjects evaluation instruments;
  \item \textbf{IT} --- trust zones and degraded-mode procedures (technical appendix);
  \item \textbf{Parent or staff representatives} --- where required by local law.
\end{itemize}

\subsection{Stakeholder next steps}

\begin{table}[htbp]
  \centering
  \caption{Immediate actions.}
  \label{tab:exec-next-steps}
  \small
  \begin{tabular}{@{}p{2.6cm}p{8.4cm}@{}}
    \toprule
    Audience & Action \\
    \midrule
    Leadership & Steering committee; approve Executive Brief circulation \\
    DPO & Open DPIA; review Box~\ref{box:system-boundaries} and checklist in full requirements doc \\
    Safety & Schedule protocol-mapping workshop and tabletops \\
    Ethics & Review transparency plan and evaluation instruments \\
    IT & Read technical appendix; plan trust zones \\
    \bottomrule
  \end{tabular}
\end{table}

\vspace{0.5em}
\noindent\textbf{Full requirements document} (\texttt{main.pdf}) includes deployment specifications,
DPIA checklist, SOPs, budgets, and governance templates.
